% 热木星（综述）
% license CCBYSA3
% type Wiki

本文根据 CC-BY-SA 协议转载翻译自维基百科\href{https://en.wikipedia.org/wiki/Hot_Jupiter}{相关文章}。

\begin{figure}[ht]
\centering
\includegraphics[width=6cm]{./figures/d9160f8109011126.png}
\caption{一幅艺术家绘制的近距离绕其母星运行的热木星示意图} \label{fig_HoJu_1}
\end{figure}
热木星（有时也称为热土星）是一类气态巨行星系外行星，被认为在物理性质上与木星相似（即木星类行星），但其轨道周期极短（(P<10) 天）$^{[1]}$。由于它们与母星距离极近、表面与大气温度很高，因而获得了“热木星”这一非正式名称$^{[2]}$。

在所有系外行星中，热木星是最容易通过径向速度法探测到的类型之一，这是因为它们在母星运动中引起的摆动幅度相对较大、变化也更快，明显强于其他已知类型的行星。其中最著名的热木星之一是 51 Pegasi b。它于 1995 年被发现，是第一颗被证实绕类太阳恒星运行的系外行星。51 Pegasi b 的轨道周期约为 4 天$^{[3]}$。
\subsection{一般特征}
\begin{figure}[ht]
\centering
\includegraphics[width=6cm]{./figures/e311dba254552faa.png}
\caption{截至 2014 年 1 月 2 日已发现的热木星（分布在左侧边缘，其中包括大多数通过凌日法探测到的行星，以黑点表示）} \label{fig_HoJu_2}
\end{figure}
\begin{figure}[ht]
\centering
\includegraphics[width=6cm]{./figures/c9637a190eec93c9.png}
\caption{具有隐藏水分子的热木星$^{[4]+}$} \label{fig_HoJu_3}
\end{figure}
尽管热木星之间存在多样性，但它们确实共享一些共同特征。
\begin{itemize}
\item 它们的定义性特征是较大的质量和极短的轨道周期，质量范围约为$0.36-11.8$个木星质量，轨道周期为$1.3-111$个地球日$^{[5]}$。其质量不能超过约 (13.6) 个木星质量，否则行星内部的压力和温度将高到足以引发氘核聚变，从而使该天体成为一颗褐矮星$^{[6]}$。
\item 大多数热木星具有近圆形轨道（低偏心率）。一般认为，其轨道会在邻近恒星的摄动或潮汐力作用下被逐渐圆化$^{[7]}$。它们是否能够在这种圆形轨道上长期稳定存在，或最终与其宿主恒星相撞，则取决于其轨道演化与物理演化之间的耦合关系，而这种耦合又通过能量耗散与潮汐形变联系在一起$^{[8]}$。
\item 许多热木星表现出异常低的平均密度。目前测得密度最低的是 TrES-4b，仅为$0.222,\mathrm{g/cm^3}$,$^{[9]}$。热木星的巨大半径尚未被完全理解，但通常认为其膨胀的大气包层可能源于强烈的恒星辐照、较高的大气不透明度、潜在的内部能量来源，以及与恒星距离极近导致行星外层超过洛希极限并被进一步拉伸$^{[9][10]}$。
\item 热木星通常处于潮汐锁定状态，即始终以同一面朝向其宿主恒星$^{[11]}$。
\item 由于其短轨道周期、相对较长的昼夜周期以及潮汐锁定特性，它们极可能拥有极端而奇特的大气环境$^{[3]}$。大气动力学模型预测，其大气将呈现显著的垂直分层结构，并伴随由辐射强迫以及热量与动量传输所驱动的强烈风场和超旋转的赤道急流$^{[12][13]}$。最新模型还预测存在多种风暴（涡旋），能够混合其大气并在不同区域之间输送冷热气体$^{[14]}$。
\item 在光球层尺度上，其昼夜温差被预测相当显著，例如基于 HD 209458 b 的模型给出的温差约为$500,\mathrm{K}$（约 $500,^\circ\mathrm{C}$，$900,^\circ\mathrm{F}$）$^{[13]}$。
\item 从统计上看，热木星更常见于 F 型和 G 型恒星周围，在 K 型恒星周围则较少，而围绕红矮星的热木星极为罕见$^{[15]}$。在对其分布进行概括时必须考虑各种观测偏差，但总体而言，其出现频率会随着恒星绝对星等的增加而呈指数式下降$^{[16]}$。
\end{itemize}
\subsection{形成与演化}

关于热木星可能起源的机制，目前主要存在三种观点。一种可能性认为，热木星是在其当前观测到的位置**原地形成**的；另一种可能性认为，它们最初在更远离恒星的位置形成，随后向内迁移；第三种观点则结合了多种动力学过程。行星向内迁移可能发生在太阳星云阶段，与气体和尘埃的相互作用有关；也可能源于与另一颗大型天体的近距离遭遇，从而使类木星行星的轨道发生不稳定$^{[3][17][18]}$。
\subsubsection{迁移}
在迁移假说中，热木星最初在**霜线**之外形成，通过行星形成的**核吸积模型**由岩石、冰和气体构建而成。随后，该行星逐渐向宿主恒星方向迁移，最终在近恒星区域形成稳定轨道$^{[19][20]}$。这种迁移可能是通过**II 型轨道迁移**平滑完成的$^{[21][22]}$；也可能是在与另一颗大质量行星发生引力散射后，被抛入高偏心率轨道，随后在与恒星的潮汐相互作用下，轨道逐渐圆化并收缩。

此外，热木星的轨道还可能受到科扎伊机制的影响，即轨道倾角与偏心率之间发生交换，使行星获得高偏心率、近日点极小的轨道，并在潮汐摩擦作用下进一步演化。这一机制通常需要一个位于更远处、且轨道倾角较大的大质量天体（另一颗行星或伴星）。观测显示，约有 (50\%) 的热木星系统中存在质量相当于或大于木星的远距离伴体，这可能导致热木星的轨道相对于恒星自转轴产生明显倾斜$^{[23]}$。

II 型迁移发生在太阳星云尚未完全消散的阶段，即系统中仍存在大量气体之时。在这一时期，高能恒星辐射和强烈的恒星风会逐渐清除残余星云$^{[24]}$。而通过其他机制（如引力散射或科扎伊机制）引发的迁移，则可能发生在气体盘消失之后。
\subsubsection{原地形成}
另一种假说认为，热木星并非由向内迁移的气态巨行星演化而来，而是其核心最初类似于常见的超级地球，在当前位置吸积气体包层并直接成长为气态巨行星，即原地形成。在这一模型中，提供核心的超级地球可能是在原地形成的，也可能是在更远处形成并在获得气体包层之前完成了迁移。

由于超级地球往往存在行星伴体，因此原地形成的热木星也可能伴随其他行星。局部生长的热木星质量增加，会对邻近行星产生多种动力学影响。例如，若热木星的轨道偏心率大于 (0.01)，扫掠的长期共振可能会增大伴星的偏心率，使其最终与热木星发生碰撞，从而导致热木星拥有异常巨大的核心；若热木星的偏心率始终较小，这种长期共振则可能使伴星轨道发生倾斜$^{[25]}$。

传统观点通常不支持原地形成模型，因为构建形成热木星所需的大质量核心，要求固体物质的面密度达到$\approx10^4\mathrm{g/cm^2}$或更高$^{[26][27][28]}$。然而，近期的观测调查发现，行星系统的内区往往富集大量超级地球型行星$^{[29][30]}$。如果这些超级地球最初在更远处形成并向内迁移，那么热木星的“原地形成”实际上也并非完全意义上的原地过程。
\subsubsection{大气流失}
如果热木星的大气层因流体动力学逃逸而被逐渐剥离，其残余核心可能演化为一种被称为冥府行星（chthonian planet）的天体。大气损失的程度取决于行星的尺寸、包层气体的成分、与恒星的轨道距离以及恒星的光度。

在一个典型系统中，一颗位于$0.02,\mathrm{AU}$处绕恒星运行的气态巨行星，在其演化寿命内可能损失约 (5\%-7\%) 的质量；若轨道距离小于 $0.015,\mathrm{AU}$，则可能发生更为剧烈的蒸发，损失行星质量中的相当一部分$^{[31]}$。迄今为止，这类天体尚未被观测到，仍停留在理论假设阶段。
\begin{figure}[ht]
\centering
\includegraphics[width=8cm]{./figures/26bf702636427d85.png}
\caption{“热木星”系外行星对比（艺术概念图）。从左上到右下依次为：WASP-12b、Boinayel、WASP-31b、Bocaprins、HD 189733b、Puli、Ditsö̀、Banksia、HAT-P-1b 和 HD 209458b。} \label{fig_HoJu_4}
\end{figure}
\subsection{具有热木星的行星系统中的类地行星}
数值模拟表明，一颗木星质量行星穿过内侧原行星盘（距恒星约 5–0.1 AU 的区域）向内迁移，其破坏性并没有此前预想的那么大。该区域中超过 60\% 的固体盘物质（包括微行星和原行星）会被向外散射，从而使行星形成盘能够在气体巨行星经过之后重新形成$^{[32]}$。在这些模拟中，当热木星通过并最终在 0.1 AU 处稳定下来之后，质量可达 2 个地球质量的行星仍然能够在宜居带内形成。由于内行星系统物质与霜线之外外行星系统物质发生混合，模拟结果显示，在热木星迁移之后形成的类地行星往往会富含水$^{[32]}$。根据一项 2011 年的研究，热木星在向内迁移过程中可能会被部分破坏，这或许可以解释在距宿主恒星 0.2 AU 以内大量存在“炽热”的地球大小至海王星大小行星的现象$^{[33]}$。

这类系统的一个例子是 WASP-47 系统。该系统包含三颗内侧行星以及一颗位于宜居带内的外侧气体巨行星。最内侧的行星 WASP-47e 是一颗质量为 6.83 个地球质量、半径为 1.8 个地球半径的大型类地行星；热木星 WASP-47b 的质量略大于木星，但半径约为 12.63 个地球半径；此外还有一颗“热海王星”WASP-47c，其质量为 15.2 个地球质量、半径为 3.6 个地球半径$^{[34]}$。开普勒-30（Kepler-30）系统也呈现出类似的轨道结构$^{[35]}$。
\subsection{失配轨道}
若干热木星（例如 HD 80606 b）的轨道与其宿主恒星的自转并不对齐，其中甚至包括逆行轨道的例子，如 HAT-P-14b$^{[36][37][38][39]}$。这种失配可能与热木星所绕行恒星光球层的温度有关。对此已有多种假说。一种假说基于潮汐耗散，认为形成热木星的机制是单一的，但会产生一系列不同的自转倾角：较冷恒星具有更强的潮汐耗散，能够抑制自转倾角（因此围绕冷恒星运行的热木星往往轨道对齐），而较热恒星则难以抑制倾角，从而导致观测到的轨道失配$^{[5]}$。另一种假说认为，变化的是宿主恒星在演化早期的自转状态，而非行星轨道本身$^{[40]}$。还有一种观点认为，热木星更可能形成于致密星团环境，在那里扰动更为频繁，邻近恒星捕获行星的引力过程也更容易发生$^{[41]}$。
\subsection{超热木星}
超热木星是指白昼侧温度高于 2200 K（约 1930 °C）的热木星。在这种白昼侧大气中，大多数分子会解离成其组成原子，并被输运到夜侧，在夜侧重新结合成分子$^{[42][43]}$。

一个例子是 TOI-1431b，于 2021 年 4 月由南昆士兰大学宣布发现。它的轨道周期仅约 2.5 天，其白昼侧温度高达 2700 K（约 2430 °C），比银河系中约 40\% 的恒星还要炽热$^{[44]}$；其夜侧温度约为 2600 K（约 2330 °C）$^{[45]}$。
\subsection{超短周期行星}
超短周期行星（USP）是一类轨道周期短于 1 天的行星，只出现在质量小于约 1.25 个太阳质量的恒星周围$^{[46][47]}$。

已确认的、轨道周期短于 1 天的凌星型热木星包括 WASP-18b、WASP-19b、WASP-43b、WASP-103b、TOI-1937A b$^{[48]}$ 以及 TOI-2109b$^{[49]}$。

截至 2025 年，已知共有 12 颗木星型超短周期行星$^{[49]}$。
\subsection{蓬松行星}
半径巨大而平均密度极低的气体巨行星有时被称为“蓬松行星”$^{[50]}$，或称“热土星”，因为它们的密度与土星相近。蓬松行星通常绕恒星近距离运行，来自恒星的强烈加热与行星内部的加热共同作用，使其大气显著膨胀。通过凌星法已探测到六颗“大半径—低密度”的行星，按发现顺序分别为：HAT-P-1b$^{[51][52]}$、CoRoT-1b、TrES-4b、WASP-12b、WASP-17b 和 Kepler-7b。一些通过径向速度法发现的热木星也可能属于蓬松行星。这类行星大多具有不高于木星的质量，因为更高质量的行星拥有更强的引力，能将自身半径维持在接近木星的尺度。事实上，质量低于木星且温度高于 1800 K 的热木星会发生极度膨胀，处在不稳定的演化路径上，最终可能导致洛希瓣溢出，并引发大气蒸发与流失$^{[53]}$。

即便考虑到来自恒星的表面加热，许多凌星型热木星的半径仍明显大于理论预期。这可能源于大气风与行星磁层之间的相互作用，在行星内部产生电流并加热行星，从而导致其膨胀。行星越热，大气电离程度越高，相互作用越强，电流越大，进而引发更显著的加热与膨胀。该理论与观测到的“行星温度与半径膨胀程度相关”的事实相一致$^{[53]}$。
\subsection{卫星}
理论研究表明，热木星不太可能拥有稳定的卫星系统。这一方面是由于其希尔球较小，另一方面是因为其所绕恒星施加的潮汐力会使任何卫星的轨道不稳定，且卫星质量越大，这种不稳定效应越强。因此，对多数热木星而言，能够长期存在的卫星往往只能是小到小行星尺度的天体$^{[54]}$。此外，热木星自身的物理演化也会决定其卫星的最终命运：可能将其“卡”在半渐近的半长轴上，或将其抛射出系统，使其经历其他未知过程$^{[55]}$。尽管如此，对 WASP-12b 的观测表明，它可能至少拥有一颗较大的系外卫星$^{[56]}$。
\subsection{绕红巨星运行的热木星}
有人提出，若气体巨行星以类似木星—太阳距离的轨道绕红巨星运行，由于将承受来自恒星的强烈辐照，它们可能表现为热木星。可以预期，在太阳演化为红巨星之后，太阳系中的木星也极有可能转变为一颗热木星$^{[57]}$。近年来发现的一些绕红巨星运行、密度异常低的气体巨行星为这一假说提供了支持$^{[58]}$。

与绕主序星运行的热木星相比，绕红巨星运行的热木星在多个方面会有所不同，尤其是它们可能会从恒星风中吸积物质；并且在假设其自转较快（未被潮汐锁定）的情况下，其热量分布会更加均匀，形成多条狭窄带状的高速喷流。由于相对于其宿主红巨星而言行星尺寸极小，加之一次凌星和一次掩食所需时间极长（可能长达数月甚至数年），利用凌星法探测这类行星将更加困难$^{[57]}$。
\subsection{恒星—行星相互作用}
自 2000 年以来的理论研究曾提出，“热木星”可能会由于其与宿主恒星之间的磁场相互作用，或二者之间的潮汐力作用，而引发恒星耀斑活动的增强。这类效应被称为“恒星—行星相互作用”（star–planet interactions，简称 SPI）。HD 189733 系统曾是被认为最可能存在此类效应、且研究最为深入的系外行星系统。

2008 年，一组天文学家首次描述称：当绕 HD 189733 A 运行的系外行星到达其轨道上的某一特定位置时，会引发恒星耀斑活动的增强。2010 年，另一研究团队发现，每当他们在该行星轨道上的某一位置观测到这颗系外行星时，都会同时探测到 X 射线耀斑。2019 年，天文学家综合分析了来自阿雷西博天文台、MOST 卫星以及自动光电望远镜（Automated Photoelectric Telescope）的数据，并结合该恒星在射电、光学、紫外和 X 射线波段的历史观测记录，对上述主张进行了系统检验。结果表明，先前的结论被夸大了：宿主恒星并未表现出与恒星耀斑或类似太阳活动区（包括黑子）相关的多种亮度和光谱特征。统计分析还发现，无论系外行星位于轨道上的何处，恒星耀斑都会出现，因此推翻了早先关于行星位置与耀斑活动相关性的说法。研究表明，宿主恒星与系外行星的磁场并不存在显著相互作用，该系统如今已不再被认为存在所谓的“恒星—行星相互作用”$^{[59]}$。此外，还有研究者曾提出，HD 189733 可能会以类似于 T Tauri 型年轻原恒星系统的速率，从其所绕行的系外行星中吸积（拉拽）物质。但后续分析显示，从这颗“热木星”伴星中被吸积的气体即便存在，其量也极其有限$^{[60]}$。
\subsection{参见}

* **大迁移假说（Grand tack hypothesis）**：一种模型，认为木星在太阳系早期曾向内迁移，随后由于与土星的轨道共振而被拉回至现今位置
* **热海王星（Hot Neptune）**：绕恒星近距离运行、大小与海王星或天王星相当的行星
* **木星类比行星（Jupiter analogue）**：与木星相似的系外行星
* **行星列表（Lists of planets）**
* **行星迁移（Planetary migration）**
* **次褐矮星（Sub-brown dwarf）**：以类似恒星形成方式产生、但质量处于行星尺度的褐矮星
\subsection{延伸阅读}

* *Origins of Hot Jupiters*，Rebekah I. Dawson，John Asher Johnson，2018 年 1 月 18 日
